\documentclass[12pt, twoside]{article}
\input{../preamble}

\fancyhead[LE]{\thepage}
\fancyhead[RO]{Markscheme}
\fancyhead[LO]{La Scuola d'Italia / Dr.\ Huson / IB Math 11 \\* 14 September 2026}

\begin{document}
\subsubsection*{1.2 Classwork: Arithmetic Sequences with Variables --- Markscheme}

\begin{enumerate}[itemsep=0.3cm]

%--- Problem 1: Review -- negative d; u_n, u_20, solve u_n = -13 [7 marks] ---
\item[1.] [Maximum mark: 7]

The sequence $83,\; 77,\; 71,\; 65,\; \ldots$

\medskip

\textbf{(a)} $d = 77 - 83 = -6$ \hfill \textbf{A1}

\emph{The negative sign is required: $d = 6$ earns no mark.}

\medskip

\textbf{(b)} $u_{n} = u_{1} + (n-1)d = 83 + (n-1)(-6)$ \hfill \textbf{M1}

$u_{n} = 83 - 6(n-1)$ \hfill \textbf{A1}

\emph{Accept the simplified form $u_{n} = 89 - 6n$, or any other equivalent
expression, for full marks. The substituted form is preferred because it shows
where the expression comes from.}

\emph{FT: award M1A1 for a correct expression formed from $u_{1} = 83$ and
the candidate's own $d$ from (a).}

\medskip

\textbf{(c)} $u_{20} = 83 - 6(20 - 1) = 83 - 114 = -31$ \hfill \textbf{A1}

\emph{FT: award A1 for the value obtained by substituting $n = 20$
into the candidate's own expression from (b).}

\medskip

\textbf{(d)} $83 - 6(n - 1) = -13$ \hfill \textbf{M1}

$-6(n - 1) = -96$

$n - 1 = 16$ \hfill \textbf{A1}

$n = 17$ \hfill \textbf{A1}

\emph{Accept the equivalent set-up $89 - 6n = -13$ leading to $6n = 102$.}

\emph{The A1 for $n - 1 = 16$ is for any correct intermediate line
that isolates $n$ (e.g.\ $6n = 102$); it does not have to be written
in this exact form.}

\emph{An equation that drops the $(n - 1)$, such as $83 - 6n = -13$
(giving $n = 16$), earns the M1 only --- unless $83 - 6n$ is the
candidate's own expression from (b), in which case FT applies.}

\emph{A candidate who states $n = 17$ from a GDC table with no
algebraic working earns the final A1 only, since the question asks
for algebraic working.}

\emph{FT: award the marks for a correct solution of the candidate's own
expression from (b) set equal to $-13$.}

\newpage

%--- Problem 2: Terms in k; equal differences -> k, the terms, u_10 [7 marks] ---
\item[2.] [Maximum mark: 7]

$u_{1} = k - 4$, \quad $u_{2} = 2k$, \quad $u_{3} = 4k - 1$.

\medskip

\textbf{(a)} $u_{2} - u_{1} = 2k - (k - 4) = k + 4$ \hfill \textbf{A1}

$u_{3} - u_{2} = (4k - 1) - 2k = 2k - 1$ \hfill \textbf{A1}

\emph{A1 for each expression, awarded independently. Accept a correct
unsimplified form, e.g.\ $2k - k + 4$.}

\emph{A sign error when subtracting the bracket, e.g.\ $2k - k - 4 = k - 4$,
earns no mark for that expression.}

\medskip

\textbf{(b)} The sequence is arithmetic, so the two differences are equal.

$k + 4 = 2k - 1$ \hfill \textbf{M1}

$k = 5$ \hfill \textbf{A1}

\emph{M1 is for any valid use of the equal-difference property: equating
the two differences, using $2u_{2} = u_{1} + u_{3}$ (so $4k = 5k - 5$), or
substituting a value of $k$ and showing that the two differences are then
equal.}

\emph{FT: award M1A1 for correctly solving the equation formed from the
candidate's own expressions in (a).}

\medskip

\textbf{(c)} $u_{1} = 1$, \quad $u_{2} = 10$, \quad $u_{3} = 19$ \hfill \textbf{A1}

\emph{All three terms required for the mark.}

\emph{FT: award A1 for the three terms obtained by substituting the
candidate's own $k$ from (b).}

\medskip

\textbf{(d)} $d = 10 - 1 = 9$

$u_{10} = 1 + (10 - 1)(9)$ \hfill \textbf{M1}

$u_{10} = 82$ \hfill \textbf{A1}

\emph{M1 is for any correct use of $u_{n} = u_{1} + (n-1)d$ with the
candidate's first term and common difference, or for continuing the list
of terms to the $10^{\text{th}}$ term.}

\emph{Using $10$ in place of $(n - 1)$, giving $1 + 10(9) = 91$, earns
the M1 only.}

\emph{FT: award M1A1 for the $10^{\text{th}}$ term correctly obtained
from the candidate's own terms in (c).}

\end{enumerate}
\end{document}
